\documentclass[11pt,a4paper]{article}
\usepackage[margin=25mm]{geometry}
\usepackage{amsmath,amssymb,amsthm,booktabs}
\newtheorem{theorem}{Theorem}
\newcommand{\Alt}{\operatorname{Alt}}
\newcommand{\Sym}{\operatorname{Sym}}
\title{The smaller group censuses in the regular-polytope algorithm paper}
\author{Human Compression project: L071 recovery}
\date{Working application; complete validation and distinct review pending}
\begin{document}\maketitle
Leemans and Mixer, \emph{Algorithms for classifying regular polytopes
with a fixed automorphism group} (2012), give twelve numerical comparisons
on physical page~4 and verify the earlier sporadic classifications on
pages~9--10. We recover these finite classifications from the acquired
theorem on enumeration of marked groups and its implementation by classes
in centralizers. The alternating groups $A_{11},\ldots,A_{14}$, the
complete $\mathrm{Co}_3$ tables and the correction to the compatible-action
criterion are treated in separate components of this same recovery.

An ordered class means a class of tuples under the full automorphism
group of the generated group. Reversal is not part of that equivalence.
Write $c$ for the number of ordered classes, $s$ for the number fixed by
reversal, and $d=(c+s)/2$ for the number after identifying duals. The
standard correspondence with regular polytopes is retained throughout.

\begin{theorem}\label{bench}
For the ten groups in Table~\ref{benchmark}, the numbers of regular
polytopes are the displayed $c,s,d$. In the twelve-group comparison of
the source, the other two entries are as follows. The group $\mathrm{Sz}(8)$
has seven classes up to duality, all of rank three, and none is self-dual.
The group $A_5\times C_2$ has eleven classes up to duality: eight of rank
three with connected diagrams and three of rank four with disconnected
diagrams. Thus its printed count eight requires the connected-diagram
restriction.
\end{theorem}
\begin{table}[ht]
\centering
\begin{tabular}{lrrr}
\toprule Group & $c$ & $s$ & $d$\\\midrule
$A_5$ & 3 & 1 & 2\\
$\mathrm{P}\Gamma\mathrm L_2(9)$ & 21 & 3 & 12\\
$S_7$ & 81 & 7 & 44\\
$\mathrm{PSL}_2(25)$ & 30 & 4 & 17\\
$\mathrm{P}\Sigma\mathrm U_3(3)$ & 60 & 2 & 31\\
$\mathrm{PGL}_2(27)$ & 190 & 6 & 98\\
$M_{12}$ & 67 & 7 & 37\\
$J_1$ & 300 & 0 & 150\\
$A_9$ & 84 & 10 & 47\\
$S_9$ & 352 & 12 & 182\\
\bottomrule
\end{tabular}
\caption{Complete classes for the benchmark groups. The last column is
the equivalence convention of the source's first comparison table.}
\label{benchmark}
\end{table}

\begin{theorem}\label{sporadic}
For the eleven sporadic simple groups of order at most $|\mathrm{He}|$,
the ordered class counts by rank are those in Table~\ref{sporadic-table}.
All other ranks have zero classes. The same statement, with the separately
displayed counts, holds for the Tits group $\,{}^2F_4(2)'$.
The complete lists supplied by Hartley and Hulpke are in bijection with
these independently enumerated ordered classes, through explicit
isomorphisms and conjugations of their actual generators.
\end{theorem}
\begin{table}[ht]
\centering
\begin{tabular}{lrrrr}
\toprule Group & Rank 3 & Rank 4 & Rank 5 & Total\\\midrule
$M_{11}$ & 0 & 0 & 0 & 0\\
$M_{12}$ & 40 & 27 & 0 & 67\\
$J_1$ & 296 & 4 & 0 & 300\\
$M_{22}$ & 0 & 0 & 0 & 0\\
$J_2$ & 261 & 31 & 0 & 292\\
$M_{23}$ & 0 & 0 & 0 & 0\\
$\mathrm{HS}$ & 465 & 111 & 4 & 580\\
$M_{24}$ & 946 & 310 & 4 & 1260\\
$\mathrm{McL}$ & 0 & 0 & 0 & 0\\
$\mathrm{He}$ & 2292 & 145 & 0 & 2437\\
$J_3$ & 584 & 2 & 0 & 586\\\midrule
$\,{}^2F_4(2)'$ & 468 & 57 & 0 & 525\\
\bottomrule
\end{tabular}
\caption{Ordered classes. The Tits group is a separate additional
classification, not one of the twenty-six sporadic simple groups.}
\label{sporadic-table}
\end{table}

\begin{proof}[Common proof of the finite census statements]
For each concrete permutation model, construct a normalizing overgroup
$A$ of $G$. The retained executions check the expected order of $G$,
normality in $A$, and
\[
                |A:C_A(G)|=|\operatorname{Aut}(G)|.
\]
Thus the conjugation action really supplies every automorphism, as
required by the acquired theorem; containment of two unrelated group
embeddings is not being used. The identifications of the concrete
library models with the named groups remain explicit standard inputs.
For $M_{12}$ we use degree $144$ to implement the outer automorphism.
For $\mathrm{McL}$ the group is the derived subgroup of the chosen
automorphism overgroup, rather than a separately stored subgroup in
unrelated coordinates. The $J_3$ overgroup is supplied by the retained
Atlas generators in degree $6156$.

The acquired initial-segment theorem enumerates one representative of
each orbit of possible next marks under the pointwise stabilizer of
the current tuple. Its intersection test is applied only to extensions
of prefixes already known to have the intersection property. The
complete theorem for classes in centralizers proves that the concrete
double-coset construction supplies exactly these candidate orbits.
Its necessary-generation test uses the full normal closure of the
eligible involution classes, as that proof requires.

Every new independent mark at least doubles the generated subgroup.
Consequently a string C-marking of $G$ has rank at most
$\lfloor\log_2|G|\rfloor$. The search computes this bound from the
group order and treats every rank through that bound. It reads no
source rank cutoff or expected count. Each of the twenty model runs
terminates normally with the complete-exhaustion record and unchanged
inputs. Their union contains $6868$ ordered tuples. The acquired
completeness theorem therefore applies to the actual complete outputs,
including the four groups for which every output list is empty.

The retained validation checks every actual tuple: each mark is an
involution in $G$, the marks generate $G$, the string commutations hold,
and the intersection equality holds for every pair of subsets of the
marking. It computes the orders of every adjacent product, giving the
complete ordered type distributions. This validation supplements the
enumeration theorem; it does not infer exhaustion from the number of
surviving tuples. Reversal is tested in the full automorphism action,
which yields $s$ and hence $d$.

The two reused benchmark rows require no further search. The acquired
complete $H_3$ census applies because $H_3\cong A_5\times C_2$ and gives
exactly the eleven stated classes. The acquired Suzuki counting theorem
gives
\[
 \frac{1}{2\cdot3}\sum_{e\mid3}\mu(3/e)(2^e-2)(2^e-1)=7.
\]
Its complete rank theorem excludes every other rank and its duality
argument excludes self-dual classes. Both entire acquired proofs and
their qualifications remain dependencies of this application.

Finally the source comparison uses all twelve original auxiliary lists,
containing $6047$ ordered markings. An isomorphism maps each source group
to our model, and a recorded automorphism then maps each source tuple
to one specified enumerated tuple. The terminal tuples are distinct,
and the two lists have equal cardinality for every group. The acquired
comparison theorem therefore gives a bijection of marked classes.
This proves the assertion about the original lists themselves, including
their ordered types; equality of the three rank totals alone would not
prove that assertion.
\end{proof}

\paragraph{Source conventions and attribution.}
The eleven ordinary sporadic groups are the literal scope of the source's
pages~9--10. Its broader description on page~2 refers to the complete
Hartley--Hulpke classifications, whose auxiliary archive also contains
the Tits group. We include the latter explicitly under that broader
reading. The twelve benchmark entries contain two of the ordinary
sporadic groups, $M_{12}$ and $J_1$; they are computed and counted only
once in the twenty-model execution. The two reused groups complete
the twenty-two distinct groups in this component.

The lists of Hartley and Hulpke and the numerical results of Leemans
and Mixer retain their original attribution. Our contribution is the
application of the common enumeration theorem, complete independent
executions, and explicit correspondence with the original generators.
Historical CPU times and memory use are measurements of the authors'
implementations and hardware; this proof does not reproduce them or
claim a performance comparison. The source's earlier $A_{10}$ work and
its open-question remark are outside the new finite classifications
proved here. No open-conjecture answer or complete Lean formalization
is claimed by these computations.
\end{document}
